\section*{अध्याय \ref{ch:b3:facies-models} --- संलक्षणी मॉडल}

\begin{solution}{exo:b3:facies-models:1}
समुद्र से स्थल की ओर: अपतट, तूफ़ानी रेत के साथ अपतट-संक्रमण,
\omterm{def:b3:facies-models:shoreface}{तटमुख}, पुलिन, दलदल वाला तटीय मैदान। प्रत्येक परिवेश अपने समुद्र-पक्षीय
पड़ोसी के ऊपर जमा हुआ: तटरेखा समुद्र की ओर खिसकी, यह एक रिग्रेशन है।
ट्रांसग्रेशन इन्हें उल्टे क्रम में ढेर करता, नीचे तटीय मैदान और ऊपर
अपतट पंक, और अक्सर वहाँ एक अपरदन सतह के साथ जहाँ
\omterm{def:b3:facies-models:shoreface}{तटमुख} गुज़रा।
\end{solution}

\begin{solution}{exo:b3:facies-models:2}
(क) वाहिका-भराव या \omterm{def:b3:facies-models:rivers}{बिंदु-पुलिन}: आधार पर कंकड़ीला अपरदन-गर्त, तिर्यक संस्तर,
फिर लहरिकाएँ और जड़ों वाला पंक। (ख) प्रगामी \omterm{def:b3:facies-models:shoreface}{तटमुख} या डेल्टा-अग्र:
नीचे तूफ़ान-रेत संस्तरों और \omterm{def:b3:facies-models:hcs}{टीलेदार तिर्यक स्तरण} वाला पंक,
ऊपर तिर्यक-संस्तरित तटमुख-रेत और पुलिन। (ग) एकत्रित गुंफित-नदी रेत, पवन-टिब्बा
या कार्बोनेट शोल: समान कण-आकार वाले तिर्यक संस्तर-समूह,
पवन-स्थिति के लिए बड़े टिब्बा-अग्र-स्तर, या शोल के लिए कंकाली और ओलाइटिक
कण।
\end{solution}

\begin{solution}{exo:b3:facies-models:3}
$\phi = (2.65 - 2.40)/(2.65 - 1.0) = 0.152$; खारे जल के साथ,
$0.25/1.55 = 0.161$।
\end{solution}

\begin{solution}{exo:b3:facies-models:4}
$e^{-2\pi z/\lambda} = e^{-\pi} = 0.043$, जब $z = \lambda/2$ होता है, लगभग
\qty{38}{m} और \qty{55}{m} गहराई पर। \qty{20}{m} पर: $e^{-2\pi\times20/75} = 0.19$
और $e^{-2\pi\times20/110} = 0.32$ सतह-गति का।
\end{solution}

\begin{solution}{exo:b3:facies-models:5}
$(60 - 47.6)/(189 - 47.6) = 0.088$; $(60 - 52.6)/(189 - 52.6) = 0.054$।
\qty{5}{\micro\second/ft} तक अनिश्चित मैट्रिक्स-मान छिद्रता को
एक-तिहाई बदल देता है: समय-औसत समीकरण को कोर पर जाँचे गए मैट्रिक्स-मान की आवश्यकता है।
\end{solution}

\begin{solution}{exo:b3:facies-models:6}
$(6.0 + 3.5 + 8.0 + 2.5)/40 = 20/40 = 0.50$। रंध्र-आयतन: $20\times0.22
\times10^6 = \qty{4.4e6}{m^3}$ प्रति वर्ग किलोमीटर।
\end{solution}

\begin{solution}{exo:b3:facies-models:7}
कैनियन के पास: केवल Ta, ऊपरी खंड अगली धारा द्वारा हटा दिए गए; जलमग्न पंखे
की ओर नीचे: Tc (लहरिकाएँ) और Td (गाद) के नीचे Te (पंक)। धारा दूरी के साथ
अपना वेग और अपना स्थूल भार खोती है; केवल महीन अंश यात्रा करता है
\qty{30}{km}, और वह एक धीमी धारा द्वारा निक्षेपित होता है जो लहरिकाएँ बनाती है,
न कि सघन निलंबन से गिरकर।
\end{solution}

\begin{solution}{exo:b3:facies-models:8}
\omterm{def:b3:facies-models:rivers}{विसर्पी नदी} का \omterm{def:b3:facies-models:rivers}{बिंदु-पुलिन}: अपरदित वाहिका-तल पर अवशेष-परत,
गहरी वाहिका की तिर्यक-संस्तरित रेत, पुलिन के शीर्ष की लहरिकायुक्त रेत,
बाढ़ के मैदान का जड़ों वाला पंक। यह अनुक्रम वाहिका को तल से
किनारे तक दर्ज करता है, अतः वाहिका लगभग \qty{9}{m} गहरी थी (पंक के संहनन के बाद
थोड़ी कम)। आधार के आर-पार नहीं: यह एक अपरदन सतह है, और
उसके नीचे जो भी था वह कोई पुराना निक्षेप हो सकता है।
\end{solution}

\begin{solution}{exo:b3:facies-models:9}
$60/40\,000 = 0.0015$, \ang{0.09}: रैंप। $300/1500 = 0.20$, \ang{11.3},
\qty{5}{m} गहरे समतल शीर्ष के बाद: \omterm{def:b3:facies-models:ramp}{कगारयुक्त मग्नतट}, किनारे पर भित्ति या शोल-कण
के साथ, और नीचे ढाल पर खंड तथा कण-प्रवाह।
\end{solution}

\begin{solution}{exo:b3:facies-models:10}
पट्टी को किसी स्थिर स्थान से गुज़रने में $w/v$ समय लगता है, जिस दौरान
उसके नीचे का तल सापेक्ष समुद्र-तल पर $A$ दर से उठता है: रेत $h = A\,w/v$ भरती है।
$2000\times0.2/20 = \qty{20}{m}$। $A = 0$ पर \omterm{def:b3:facies-models:shoreface}{तटमुख} केवल
अपनी ही रूपरेखा की मोटाई छोड़ता है, एक पतली परत, क्योंकि पट्टियाँ एक-दूसरे पर
ऊर्ध्व-निक्षेपण के बिना फिसलती हैं। अनुपात $A/v$ प्रक्षेप-पथ के कोण की स्पर्शज्या है
अध्याय~\ref{ch:b3:sequence-stratigraphy} की प्रतिज्ञप्ति में: प्रत्येक संलक्षणी की मोटाई
उसकी चौड़ाई गुणा उसी स्पर्शज्या के बराबर होती है।
\end{solution}

\begin{solution}{exo:b3:facies-models:11}
घनत्व: $(2.65 - 2.30)/1.65 = 0.212$; ध्वनिक: $(86 - 52.6)/136.4 = 0.245$।
$\phi = 0.245$ मानने पर, $\rho_{\mathrm{ma}} = (2.30 - 0.245)/(1 - 0.245) =
\qty{2.72}{g/cm^3}$: क्वार्ट्ज़ से घना, जैसा कैल्साइट (2.71) से सीमेंटित
बलुआ पत्थर या कैल्साइट-कण धारण करने वाली चट्टान का होता है।
\end{solution}

\begin{solution}{exo:b3:facies-models:12}
कोयले पर अपतट पंक असंलग्न परिवेशों को संपर्क में लाता है:
एक \omterm{def:b3:sequence-stratigraphy:parasequence}{प्लावन सतह}, यहाँ तरंग-अपरदन की एक \omterm{def:b3:sequence-stratigraphy:surfaces}{ट्रांसग्रेसिव सतह}, शंख-मलबा
एक अवशेष-परत। तटरेखा के स्थल की ओर खिसकने पर तरंगों ने तल को
लगभग \omterm{def:b3:facies-models:shoreface}{सामान्य-मौसम तरंग-आधार} तक काट दिया: पुलिन, ऊपरी \omterm{def:b3:facies-models:shoreface}{तटमुख} और
तटीय मैदान का शीर्ष, लगभग \qty{10}{m} तक का परिच्छेद, संभावित रूप से
लुप्त हो सकता है।
\end{solution}

\begin{solution}{pb:b3:facies-models:1}
\textbf{1.} चॉक (ठंडा, कम घनत्व), \qty{766}{m} से शेल (गर्म), चूना-पत्थर
\qty{866}{m} से (ठंडा, सघन, तेज़), बलुआ पत्थर \qty{1046}{m} से (घंटी-आकार
आकृतियाँ), शिस्ट \qty{1166}{m} से।

\textbf{2.} चॉक \qty{766}{m}, शेल 100, चूना-पत्थर 180, बलुआ पत्थर
\qty{120}{m}; आवरण का आधार $441.0 - 1166 = \qty{-725}{m}$।

\textbf{3.} $(2.71 - 2.03)/1.71 = 0.40$ सतह के निकट, $(2.71 -
2.36)/1.71 = 0.20$ \qty{750}{m} पर: चॉक अपनी आधी छिद्रता खोता है
अधिभार के भार के नीचे, संहनन द्वारा।

\textbf{4.} इसके कण अनिश्चित घनत्व वाली मृत्तिकाएँ हैं, और इसका जल
उनसे आंशिक रूप से बँधा है: समीकरण की मान्यताएँ विफल होती हैं।

\textbf{5.} क्रिस्टलीय आधारशिला: आवरण पार किया जाता है, और
नीचे किसी भंडार की अपेक्षा नहीं होती।

\textbf{6.} $(2.71 - 2.573)/1.71 = 0.080$.

\textbf{7.} $(58.9 - 47.6)/141.4 = 0.080$: वही।

\textbf{8.} $(2.65 - 2.32)/1.65 = 0.200$.

\textbf{9.} $(79.9 - 52.6)/136.4 = 0.200$: वही।

\textbf{10.} $0.33/1.55 = 0.213$, एक बिंदु और एक-तिहाई अधिक।

\textbf{11.} बलुआ पत्थर, जिसकी रेत में \qty{20}{\percent} छिद्रता है;
चूना-पत्थर के ऊपर का शेल उसके लिए आवरण चट्टान है, और बाढ़ के मैदान के पंक बलुआ पत्थर को
कक्षों में बाँटते हैं।

\textbf{12.} घंटी जैसा आकार: तिर्यक संस्तरित रेत में कम गामा-किरण मान, जो
महीन होते अंतराल से होकर मूल-युक्त पंक के ऊँचे मानों तक चढ़ता है, एक तीक्ष्ण आधार पर।

\textbf{13.} \omterm{def:b3:facies-models:association}{ऊपर की ओर महीन होता अनुक्रम}: अपरदनात्मक आधार पर वाहिका-तल का अवशेष-निक्षेप,
वाहिका के टिब्बे (तिर्यक संस्तरण), \omterm{def:b3:facies-models:rivers}{बिंदु-पुलिन} का शीर्ष (लहरिकाएँ), बाढ़ का मैदान
(मूल-युक्त पंक)।

\textbf{14.} वाहिका, उसका \omterm{def:b3:facies-models:rivers}{बिंदु-पुलिन} आंतरिक तट पर, तटबंध और एक
मृदा-युक्त बाढ़ का मैदान, साथ-साथ।

\textbf{15.} रेत और महीन होते अंतराल की मोटाई लगभग $10.4 + 3.2 =
\qty{13.6}{m}$, संहनन से पहले थोड़ी अधिक।

\textbf{16.} एक \omterm{def:b3:facies-models:rivers}{विसर्पी नदी}, जिसकी वाहिका-पट्टी द्रोणी के धँसने के दौरान
बार-बार \omterm{def:b3:facies-models:rivers}{अवल्शन} द्वारा खिसकी, हर नई वाहिका
पिछली द्वारा छोड़े गए बाढ़ के मैदान में कटाव करती हुई।

\textbf{17.} नहीं: हर आधार अपरदनात्मक है; उसके नीचे का बाढ़-मैदान का पंक और
उसके ऊपर की वाहिका-रेत का पड़ोसी होना आवश्यक नहीं था।

\textbf{18.} $(60 - 30)/(105 - 30) = 0.40$.

\textbf{19.} न्यूनतम: $14\times(0.60 + 0.40\times0.25) = \qty{9.80}{m}$;
अधिकतम: $16\times(0.65 + 0.40\times0.20) = \qty{11.68}{m}$।

\textbf{20.} शुद्ध रेत \qty{58.8}{m}; महीन होते अंतराल \qty{33.4}{m}; पंक
\qty{27.8}{m}।

\textbf{21.} शुद्ध $= 58.8 + 0.40\times33.4 = \qty{72.2}{m}$; \omterm{def:b3:facies-models:ntg}{शुद्ध-से-सकल
अनुपात} $72.2/120 = 0.60$।

\textbf{22.} भिन्न 0.20 और 0.60 पर: \qty{65.5}{m} और \qty{78.8}{m}, यानी
केवल कट-ऑफ़ के चुनाव से \qty{13}{m} का फैलाव।

\textbf{23.} $72.2\times0.20\times10^6 = \qty{1.44e7}{m^3}$ प्रति वर्ग
किलोमीटर।

\textbf{24.} वार्डन-2 पर ट्राइऐसिक बलुआ पत्थर की शुद्ध रेत-मोटाई
\textbf{\qty[text-series-to-math=true]{72.2}{m}} है।
\end{solution}
